```latex
\documentclass[12pt]{article}
\usepackage{geometry}
\usepackage{amsmath}
\usepackage{booktabs}
\usepackage{siunitx}
\usepackage{caption}
\usepackage{graphicx}
\usepackage{float}

\geometry{margin=1in}

\title{Experiment Report: Resonance Frequency Analysis in a Tensioned String System}
\author{}
\date{}

\begin{document}

\maketitle

\section*{Objective}
The objective of this experiment is to investigate the relationship between tension (\(T\)) applied to a string and its resonant frequencies, specifically the experimental frequency (\(f_e\)) and the calculated frequency (\(f_c\)). The fractional deviation \(\Delta f / f\) between these two values is analyzed to assess the accuracy of theoretical predictions.

\section*{Theory}
For a stretched string fixed at both ends, the fundamental resonant frequency is given by:
\[
f = \frac{1}{2L} \sqrt{\frac{T}{\mu}}
\]
where:
\begin{itemize}
  \item \(f\) is the resonant frequency (Hz),
  \item \(T\) is the tension in the string (N),
  \item \(L\) is the length of the vibrating segment (m),
  \item \(\mu\) is the linear mass density of the string (kg/m).
\end{itemize}

In this experiment, the theoretical frequency \(f_c\) is computed using known values of \(T\) and assumed constants for \(L\) and \(\mu\). The measured frequency \(f_e\) is obtained from experimental observations. The relative deviation \(\Delta f / f\) is defined as:
\[
\frac{\Delta f}{f} = \left| \frac{f_e - f_c}{f_c} \right|
\]
This quantity provides a measure of the discrepancy between theory and experiment.

\section*{Data and Analysis}

The following table presents the collected data:

\begin{table}[H]
\centering
\caption{Measured and Calculated Resonant Frequencies vs. Applied Tension}
\label{tab:data}
\begin{tabular}{cccccc}
\toprule
\(T\) (N) & \(f_e\) (Hz) & \(f_c\) (Hz) & \(\Delta f / f\) \\
\midrule
9.801 & 45.4 & 49.2 & 0.08 \\
19.602 & 64.3 & 67.4 & 0.05 \\
29.403 & 78.7 & 83.7 & 0.06 \\
39.204 & 90.9 & 91.9 & 0.01 \\
49.005 & 101.6 & 107.8 & 0.06 \\
\bottomrule
\end{tabular}
\end{table}

The average value of \(\Delta f / f\) is calculated as:
\[
\bar{\Delta f / f} = \frac{0.08 + 0.05 + 0.06 + 0.01 + 0.06}{5} = \frac{0.26}{5} = 0.052
\]

Thus, the average deviation is \(5.2\%\), indicating a generally good agreement between experimental and theoretical results, with the largest deviation occurring at the lowest tension.

\section*{Uncertainty Analysis}
Assuming that the primary sources of uncertainty are in the measurement of frequency (\(f_e\)) and tension (\(T\)), and given that the tension values are precise to ±0.001 N (from the provided data precision), we estimate the uncertainty in \(f_e\) to be approximately ±0.5 Hz based on typical digital frequency counter resolution.

Using error propagation for the frequency formula:
\[
\frac{\Delta f}{f} \approx \frac{1}{2} \frac{\Delta T}{T} + \frac{1}{2} \frac{\Delta \mu}{\mu}
\]
Since \(\mu\) is assumed constant and known precisely, the dominant uncertainty comes from \(T\). For example, at \(T = 9.801\,\mathrm{N}\), \(\Delta T / T = 0.001 / 9.801 \approx 0.0001\), which contributes only about 0.005\% to the uncertainty — negligible compared to observed deviations.

Therefore, the observed discrepancies are likely due to systematic effects such as:
\begin{itemize}
  \item Imperfect string uniformity,
  \item Non-ideal boundary conditions (e.g., clamping not perfectly rigid),
  \item Air resistance or damping,
  \item Measurement timing errors during resonance detection.
\end{itemize}

\section*{Discussion}
The results show that the experimental frequencies are consistently lower than the calculated ones, except at the highest tension where they are very close. This suggests a possible underestimation of the effective length or overestimation of the string's mass density in the theoretical model.

Notably, the smallest deviation (\(0.01\)) occurs at \(T = 39.204\,\mathrm{N}\), suggesting optimal conditions near this tension. The increase in deviation at low tension may indicate greater sensitivity to small imperfections in setup or measurement noise.

The overall trend of \(\Delta f / f\) shows no clear correlation with increasing tension, implying that the source of error is not purely proportional to \(T\). This supports the idea of systematic rather than random error.

\section*{Conclusion}
The experiment successfully demonstrates the dependence of resonant frequency on string tension, with experimental results showing reasonable agreement with theoretical predictions. The average fractional deviation is \(5.2\%\), which is within acceptable limits for undergraduate laboratory settings. Minor discrepancies are attributed to practical limitations such as non-ideal boundary conditions and measurement uncertainties. Future improvements could include more precise frequency measurement techniques and better control of string end constraints.

\end{document}
```